% Options for packages loaded elsewhere
\PassOptionsToPackage{unicode}{hyperref}
\PassOptionsToPackage{hyphens}{url}
\documentclass[
  11pt,
  a4paper,
]{article}
\usepackage{xcolor}
\usepackage[margin=22mm]{geometry}
\usepackage{amsmath,amssymb}
\setcounter{secnumdepth}{-\maxdimen} % remove section numbering
\usepackage{iftex}
\ifPDFTeX
  \usepackage[T1]{fontenc}
  \usepackage[utf8]{inputenc}
  \usepackage{textcomp} % provide euro and other symbols
\else % if luatex or xetex
  \usepackage{unicode-math} % this also loads fontspec
  \defaultfontfeatures{Scale=MatchLowercase}
  \defaultfontfeatures[\rmfamily]{Ligatures=TeX,Scale=1}
\fi
\usepackage{lmodern}
\ifPDFTeX\else
  % xetex/luatex font selection
\fi
% Use upquote if available, for straight quotes in verbatim environments
\IfFileExists{upquote.sty}{\usepackage{upquote}}{}
\IfFileExists{microtype.sty}{% use microtype if available
  \usepackage[]{microtype}
  \UseMicrotypeSet[protrusion]{basicmath} % disable protrusion for tt fonts
}{}
\makeatletter
\@ifundefined{KOMAClassName}{% if non-KOMA class
  \IfFileExists{parskip.sty}{%
    \usepackage{parskip}
  }{% else
    \setlength{\parindent}{0pt}
    \setlength{\parskip}{6pt plus 2pt minus 1pt}}
}{% if KOMA class
  \KOMAoptions{parskip=half}}
\makeatother
\usepackage{longtable,booktabs,array}
\newcounter{none} % for unnumbered tables
\usepackage{calc} % for calculating minipage widths
% Correct order of tables after \paragraph or \subparagraph
\usepackage{etoolbox}
\makeatletter
\patchcmd\longtable{\par}{\if@noskipsec\mbox{}\fi\par}{}{}
\makeatother
% Allow footnotes in longtable head/foot
\IfFileExists{footnotehyper.sty}{\usepackage{footnotehyper}}{\usepackage{footnote}}
\makesavenoteenv{longtable}
\setlength{\emergencystretch}{3em} % prevent overfull lines
\providecommand{\tightlist}{%
  \setlength{\itemsep}{0pt}\setlength{\parskip}{0pt}}
\usepackage{bookmark}
\IfFileExists{xurl.sty}{\usepackage{xurl}}{} % add URL line breaks if available
\urlstyle{same}
\hypersetup{
  pdftitle={Linear forms in logarithms: the problem{,} the trivial bounds{,} and effectivity},
  pdfauthor={GPT-6.1 Sol (OpenAI), Ultra},
  hidelinks,
  pdfcreator={LaTeX via pandoc}}

\title{Linear forms in logarithms: the problem, the trivial bounds, and
effectivity}
\author{GPT-6.1 Sol (OpenAI), Ultra}
\date{October 2026}

\providecommand{\Eta}{\mathrm{H}}
\begin{document}
\maketitle

\emph{Draft. Public domain (CC0).}

Two integers can be very large and still differ by one. To study a
difference such as \(2^x-3^y\), divide by one of the powers. The
question becomes how close \(2^x3^{-y}\) can be to \(1\). Taking
logarithms changes multiplication into addition, but introduces a second
question: how much cancellation can occur in \(x\log 2-y\log 3\)?

This lesson develops the elementary estimates that make the question
precise. It proves the passage between the multiplicative and additive
forms, a height bound, a box-principle construction of small forms, and
an explicit finiteness consequence of a polynomial lower bound. Its
height prerequisite is the written lesson
\href{../../TR-TRANS/TR-TRANS-02.html}{\emph{Heights of algebraic
numbers}} in \emph{Transcendental numbers}: Theorem 2.4 supplies the
normalization, Proposition 2.5 the product and sum inequalities, and
Theorem 2.8 the one-place Liouville inequality. The freely accessible
treatments {[}Waldschmidt 1992{]}, {[}Waldschmidt 2000{]} and {[}Evertse
2019{]} provide the comparisons discussed below. The logarithm methods
of Baker and later authors are credited in the references.

\subsection{1. Start with a Diophantine
equation}\label{start-with-a-diophantine-equation}

If \(2^x-3^y=k\), then

\[
2^x3^{-y}-1=k3^{-y}.
\]

For fixed nonzero \(k\), the right side is exponentially small as \(y\)
grows. A lower bound that decreases only as a power of \(\max(x,y)\)
must eventually contradict this identity. This is the central mechanism
of effective applications.

Let \(\alpha_1,\ldots,\alpha_n\) be nonzero algebraic numbers, let
\(b_1,\ldots,b_n\) be integers, and choose logarithms \(\ell_j\) with
\(e^{\ell_j}=\alpha_j\). Put

\[
P=\prod_{j=1}^n\alpha_j^{b_j},\qquad
\Delta=P-1,\qquad
\Lambda=\sum_{j=1}^n b_j\ell_j,\qquad
B=\max\{3,|b_1|,\ldots,|b_n|\}.
\]

Thus \(e^\Lambda=P\). In particular,

\[
\Delta=0\quad\Longleftrightarrow\quad\Lambda\in2\pi i\mathbb Z.
\]

If the \(\alpha_j\) are positive real numbers and the \(\ell_j\) are
real logarithms, this reduces to
\(\Delta=0\Longleftrightarrow\Lambda=0\). A family is
\emph{multiplicatively independent} if \(P=1\) forces every \(b_j=0\).
For positive real numbers this is exactly the rational linear
independence of their real logarithms.

For \(2\) and \(3\), independence follows from unique factorization:
\(2^u3^v=1\) implies \(u=v=0\). More generally, independence of positive
rationals can be checked by linear algebra on the vectors of their prime
valuations. For example, \(2,3,6\) are dependent even though no two of
them are powers of each other.

An \emph{inhomogeneous} logarithmic form has the shape

\[
\beta_0+\beta_1\ell_1+\cdots+\beta_n\ell_n,
\qquad \beta_j\in\overline{\mathbb Q}.
\]

It need not exponentiate to an algebraic number. Its exponential is
\(e^{\beta_0}\prod_j\alpha_j^{\beta_j}\), where each algebraic power
means \(e^{\beta_j\ell_j}\) for the chosen logarithm. For instance,
\(\pi+\log 2=-i\log(-1)+\log2\) is a homogeneous algebraic-coefficient
logarithmic form if \(\log(-1)=i\pi\).

\subsection{2. Which logarithm is
small?}\label{which-logarithm-is-small}

Our principal logarithm has imaginary part in \((-\pi,\pi]\). We only
use its holomorphic restriction near \(1\), so its convention on the
negative real axis creates no ambiguity in the following estimates.

\textbf{Lemma 1.1 (local comparison).} If \(|w|<1/2\), then

\[
|\operatorname{Log}(1+w)|\le2|w|.
\]

If \(|z|\le1/2\), then

\[
|e^z-1|\le2|z|.
\]

\textbf{Proof.} The line segment from \(1\) to \(1+w\) stays in the
right half-plane. Hence

\[
\operatorname{Log}(1+w)=\int_0^1\frac{w}{1+tw}\,dt.
\]

The integrand has absolute value at most \(|w|/(1-|w|)\le2|w|\).
Likewise

\[
e^z-1=z\int_0^1 e^{tz}\,dt,
\]

and \(e^{|z|}\le e^{1/2}<2\). This proves both estimates. \(\square\)

\textbf{Lemma 1.2 (branch correction).} Take principal logarithms of the
\(\alpha_j\). If \(|\Delta|<1/2\), there is an even integer \(b_0\) such
that

\[
|b_0|\le\sum_j|b_j|+1\le nB+1,
\qquad
\left|b_0\log(-1)+\sum_jb_j\log\alpha_j\right|\le2|\Delta|.
\]

\textbf{Proof.} Set \(u=\operatorname{Log}(1+\Delta)\). Since
\(e^u=e^\Lambda\), the elementary description of the kernel of the
exponential gives \(\Lambda-u=2\pi i k\) for an integer \(k\). Choose
\(b_0=-2k\). Then the corrected form equals \(u\), so Lemma 1.1 gives
its bound. Moreover,

\[
\pi|b_0|=|\operatorname{Im}(\Lambda-u)|
\le\pi\sum_j|b_j|+|\operatorname{Im}u|.
\]

The number \(1+\Delta\) has positive real part, so
\(|\operatorname{Im}u|<\pi/2\). The claimed integer bound follows.
\(\square\)

The correction cannot be omitted. If \(\alpha_1=\alpha_2=-1\) and
\(b_1=b_2=1\), then \(P=1\) but \(\Lambda=2\pi i\). There are examples
with a small nonzero \(\Delta\) too: take \(\alpha_1=-1\),
\(\alpha_2=-(1+1/N)\), and both coefficients equal to \(1\). Then
\(\Delta=1/N\), while \(\Lambda=\log(1+1/N)+2\pi i\).

For positive real bases, \(\Lambda\) is real and no correction is
needed. For complex bases, smallness of \(P-1\) means closeness of
\(\Lambda\) to the period lattice, not necessarily to its origin.

\subsection{3. What arithmetic alone
gives}\label{what-arithmetic-alone-gives}

We use the absolute logarithmic Weil height \(h\). For a degree-\(d\)
number \(\gamma\), whose primitive minimal polynomial has positive
leading coefficient \(a\) and roots \(\gamma_1,\ldots,\gamma_d\), our
normalization is

\[
h(\gamma)=\frac1d\left(\log a+\sum_{r=1}^d\log\max\{1,|\gamma_r|\}\right).
\]

The prerequisite height facts are

\[
h(uv)\le h(u)+h(v),\quad h(u^{-1})=h(u),\quad
h(u+v)\le h(u)+h(v)+\log2,
\]

and the one-place Liouville inequality: if
\(0\ne u\in K\subset\mathbb C\) and \([K:\mathbb Q]=D\), then

\[
\log|u|\ge-Dh(u).
\]

The written internal proof provider is
\href{../../TR-TRANS/TR-TRANS-02.html}{\emph{Heights of algebraic
numbers}}, Theorem 2.8, together with Theorem 2.4 and Proposition 2.5.
Its statement covers every number field and every nonzero algebraic
element, with this absolute normalization. {[}Waldschmidt 1992, Chapter
3{]} gives a freely accessible complementary treatment with the same
absolute-height convention.

\textbf{Proposition 1.3 (elementary lower bound).} If \(\Delta\ne0\) and
\(K=\mathbb Q(\alpha_1,\ldots,\alpha_n)\) has degree \(D\), then

\[
\log|\Delta|\ge-D\left(\log2+\sum_j|b_j|h(\alpha_j)\right)
\ge-D\left(\log2+B\sum_jh(\alpha_j)\right).
\]

\textbf{Proof.} Repeated multiplication and inversion give
\(h(P)\le\sum_j|b_j|h(\alpha_j)\). Since \(h(1)=0\), we have
\(h(P-1)\le h(P)+\log2\). Apply Liouville's inequality to the nonzero
element \(P-1\) of \(K\). \(\square\)

There is a sharper elementary estimate for rationals. Write
\(\alpha_j=u_j/v_j\) in lowest terms, with nonzero integers \(u_j,v_j\),
and put \(A_j=\max\{|u_j|,|v_j|\}\). Before any cancellation, write
\(P=U/V\) with integers \(U,V\ne0\). The denominator satisfies
\(|V|\le\prod_j A_j^{|b_j|}\). If \(P\ne1\), the integer \(U-V\) is
nonzero, so

\[
|\Delta|=\frac{|U-V|}{|V|}\ge\prod_jA_j^{-|b_j|}.
\]

For \(\Delta=2^m3^{-n}-1\), with \(m,n\ge0\), one can use the actual
denominator \(3^n\), obtaining \(|\Delta|\ge3^{-n}\). This improves the
uniform denominator estimate \(2^{-m}3^{-n}\), but still falls
exponentially with \(n\).

For example, \(3^5/2^8-1=-13/256\). The actual denominator gives
\(13/256\ge1/256\); the uniform product estimate gives only
\(1/(3^5 2^8)\). The problem is not that the elementary bound is false.
It allows far more cancellation than the exponents usually produce.

\subsection{4. Why no fixed positive lower bound
works}\label{why-no-fixed-positive-lower-bound-works}

\textbf{Lemma 1.4 (box principle).} Let \(n\ge2\), let
\(a_1,\ldots,a_n\ge2\) be integers, and set \(A=\max_j a_j\). For every
integer \(B>2n\log A\), there are integers \(b_j\) such that

\[
0<\max_j|b_j|\le B,
\qquad
\left|\prod_j a_j^{b_j}-1\right|\le\frac{2n\log A}{B^{n-1}}.
\]

\textbf{Proof.} Consider the \((B+1)^n\) labelled points

\[
s(c)=\sum_j c_j\log a_j,\qquad c\in\{0,\ldots,B\}^n,
\]

in an interval of length at most \(nB\log A\). If two labels give the
same point, their difference supplies a nonzero vector with product
exactly \(1\), which proves the assertion. Otherwise, sort the points.
One of the \((B+1)^n-1\) successive gaps is at most

\[
\frac{nB\log A}{(B+1)^n-1}\le\frac{n\log A}{B^{n-1}}.
\]

Subtract the corresponding labels. The resulting nonzero vector has
coefficients of absolute value at most \(B\), and its logarithmic form
has absolute value at most the displayed quantity. Since \(n\ge2\) and
\(B>2n\log A\), this is less than \(1/2\). Lemma 1.1 finishes the proof.
\(\square\)

For multiplicatively independent bases, the products constructed here
are different from \(1\). Their distances from \(1\) tend to zero as
\(B\to\infty\). The same argument works for positive rationals: the
points lie in an interval of length at most \(B\sum_j|\log\alpha_j|\).
Thus there cannot be a positive lower bound independent of the
coefficients.

Combining the elementary lower bound and the box principle leaves a wide
range:

\[
e^{-cB}\quad\hbox{versus}\quad B^{-(n-1)}.
\]

The first holds for every nonzero form; the second describes forms that
exist for arbitrarily large coefficient bounds. A theorem of the shape
\(|\Delta|\ge B^{-C}\), with a fixed \(C\), closes the most important
part of that range. It does not assert that every coefficient vector
nearly attains the bound.

\subsection{5. From a polynomial bound to a finite
search}\label{from-a-polynomial-bound-to-a-finite-search}

\textbf{Proposition 1.5.} Suppose an explicitly known \(C>0\) satisfies

\[
|2^x3^{-y}-1|\ge\max\{3,x,y\}^{-C}
\]

whenever \(x,y\ge0\) and \(2^x3^{-y}\ne1\). For a nonzero integer \(k\),
every solution of \(2^x-3^y=k\) satisfies

\[
\begin{aligned}
\max\{x,y\}&\le\max\{3,R(C,k)\},\\
R(C,k)&=\left(\frac{C+\sqrt{C^2+4\log2\,\log(2|k|)}}{2\log2}\right)^2.
\end{aligned}
\]

\textbf{Proof.} Put \(B=\max\{3,x,y\}\). The equation and the assumed
bound give

\[
3^y\le|k|B^C,
\qquad 2^x\le3^y+|k|\le2|k|B^C.
\]

If \(\max(x,y)\ge3\), then \(B=\max(x,y)\) and
\(2^B\le\max(2^x,3^y)\le2|k|B^C\). Taking logarithms yields

\[
B\log2\le\log(2|k|)+C\log B.
\]

For \(t\ge1\), \(\log t\le\sqrt t\): differentiation shows the maximum
of \(\log t/\sqrt t\) is \(2/e<1\). With \(s=\sqrt B\), the preceding
inequality implies

\[
(\log2)s^2-Cs-\log(2|k|)\le0.
\]

The positive root bounds \(s\), giving \(B\le R(C,k)\). The remaining
solutions have both exponents below \(3\). \(\square\)

This deliberately generous bound is enough to make a finite search
possible. Later reductions replace large initial bounds by much smaller
ones. If \(C\) is merely known to exist, Proposition 1.5 proves
finiteness but does not specify the search range.

\subsubsection{General differences of
powers}\label{general-differences-of-powers}

Normalizing by the larger power gives a symmetric estimate. This is the
argument of {[}Evertse 2019{]}, Corollary 5.5; here we retain a
prefactor, include zero exponents, and supply the resulting search
bound.

\textbf{Corollary (general differences of powers).} Fix integers
\(a,b\ge2\). Suppose known constants \(0<c\le1\) and \(C>0\) satisfy \[
|a^u b^v-1|\ge c\max\{3,|u|,|v|\}^{-C}
\] for every integer pair \((u,v)\) with \(a^u b^v\ne1\). For
\(m,n\ge0\) with \(a^m\ne b^n\), put \(B=\max\{3,m,n\}\). Then \[
|a^m-b^n|\ge c\max\{a^m,b^n\}B^{-C}.
\] In particular, if \(a^m-b^n=k\ne0\), put \[
\begin{aligned}
T&=\log(|k|/c),\\
R_*&=\left(\frac{C+\sqrt{C^2+4T\log2}}{2\log2}\right)^2.
\end{aligned}
\] Then \[
\max(m,n)\le\max\{2,R_*\}.
\]

\textbf{Proof.} If \(a^m>b^n\), apply the assumed estimate to
\(a^{-m}b^n\) and multiply by \(a^m\). In the other case, apply it to
\(a^m b^{-n}\) and multiply by \(b^n\). The coefficient bound is \(B\)
in either orientation. This proves the first inequality.

When \(\max(m,n)\ge3\), at least one exponent equals \(B\), so
\(\max(a^m,b^n)\ge2^B\). Consequently \[
B\log2\le\log(|k|/c)+C\log B.
\] Use \(\log B\le\sqrt B\) as in Proposition 1.5 and solve the
quadratic inequality in \(\sqrt B\). Its constant term is nonpositive,
since \(|k|\ge1\) and \(c\le1\). If both exponents are below three,
their maximum is at most two. These cases give the stated bound.
\(\square\)

The fixed-base estimate (5.47), proved later in
\href{../TR-BAKER-05.html\#10-consequences-and-later-refinements}{\emph{Effective
lower bounds II: proof of Baker's theorem}}, supplies an effective \(C\)
with \(c=1\). Its proof follows the complete effective theorem in that
lesson. Neither that estimate nor the corollary requires multiplicative
independence. They apply, for example, to \(4^m-8^n=k\), provided the
difference is nonzero. An elementary valuation argument solves one such
example in Exercise 6.

\subsection{6. Effectivity and the shape of the later
theorems}\label{effectivity-and-the-shape-of-the-later-theorems}

A constant is \emph{effectively computable from given data} if the proof
supplies a terminating procedure producing a valid value from those
data. A formula is one way to do this. A finite algorithm is another. An
ineffective existence proof supplies neither, even if its conclusion is
true.

Here is a concrete contrast. For \(n\ge2\) multiplicatively independent
positive integers \(a_1,\ldots,a_n\), Gelfond's many-logarithm result
says that, for every \(\delta>0\), there exists \(c_\delta>0\) such that

\[
|a_1^{b_1}\cdots a_n^{b_n}-1|\ge c_\delta e^{-\delta B}
\]

for nonzero coefficient vectors, where \(B=\max\{3,|b_j|\}\). The quoted
result is ineffective in \(c_\delta\). Gelfond's earlier two-logarithm
estimate gives, for each \(\varepsilon>0\), an effective
\(c_\varepsilon>0\) in a lower bound

\[
|a_1^{b_1}a_2^{b_2}-1|\ge c_\varepsilon
\exp\bigl(- (\log B)^{5+\varepsilon}\bigr).
\]

These historical attributions and the distinction between the original
effective and ineffective methods are discussed in {[}Waldschmidt
2000{]}, Section 1.2, especially Theorem 1.9, and {[}Waldschmidt
1992{]}, Introduction. Both free treatments state this historical
comparison for integer bases. The stronger fixed-base polynomial
estimate in this course applies to positive rationals as well. The
following calculation proves both weaker numerical conclusions in that
full setting with effective constants.

\textbf{Lemma 1.6 (weaker bounds from a polynomial estimate).} Suppose
\(|\Delta|\ge aB^{-C}\), with known \(a,C>0\) and \(B\ge3\). For every
\(\delta>0\) and every \(q>1\),

\[
|\Delta|\ge a\exp\bigl(C-C\log(C/\delta)\bigr)e^{-\delta B},
\]

and

\[
|\Delta|\ge a\exp\left(-(q-1)(C/q)^{q/(q-1)}\right)
\exp\bigl(- (\log B)^q\bigr).
\]

\textbf{Proof.} On \(x>0\), the function \(\delta x-C\log x\) has its
minimum at \(x=C/\delta\), where its value is \(C-C\log(C/\delta)\).
Thus \(B^{-C}e^{\delta B}\) is at least the first exponential constant.
On \(t\ge0\), the minimum of \(t^q-Ct\) occurs at \(t=(C/q)^{1/(q-1)}\)
and equals \(-(q-1)(C/q)^{q/(q-1)}\). Use \(t=\log B\) for the second
bound. \(\square\)

With \(q=5+\varepsilon\), the second formula has the two-logarithm shape
above. For fixed positive rational bases, the polynomial estimate is
proved in the effective-bound lessons, specifically (5.47), from the
full effective Baker theorem and the local comparison of Section 2.
Consequently both weaker numerical bounds admit effective constants
whenever the product is not one; the historical many-logarithm proof did
not supply them. Lemma 1.6 states exactly where that later input is
used. The elementary arguments in Sections 2--4 do not depend on it.

The following table describes shapes, not interchangeable theorem
statements. The heights, branches and nonvanishing conditions must be
specified afresh whenever a theorem is applied. Write
\(\Omega=\prod_j\log A_j\), where the admissible \(A_j\) depend on the
theorem.

{\def\LTcaptype{none} % do not increment counter
\begin{longtable}[]{@{}
  >{\raggedright\arraybackslash}p{(\linewidth - 4\tabcolsep) * \real{0.3333}}
  >{\raggedright\arraybackslash}p{(\linewidth - 4\tabcolsep) * \real{0.3333}}
  >{\raggedright\arraybackslash}p{(\linewidth - 4\tabcolsep) * \real{0.3333}}@{}}
\toprule\noalign{}
\begin{minipage}[b]{\linewidth}\raggedright
Result
\end{minipage} & \begin{minipage}[b]{\linewidth}\raggedright
Characteristic conclusion
\end{minipage} & \begin{minipage}[b]{\linewidth}\raggedright
Treatment in this course
\end{minipage} \\
\midrule\noalign{}
\endhead
\bottomrule\noalign{}
\endlastfoot
Baker's qualitative theorem & Rationally independent logarithms remain
independent over the algebraic numbers, even after adjoining \(1\) &
Proved by an auxiliary function and extrapolation \\
Baker's algebraic-coefficient measure & A nonzero form has size
\(>B^{-C}\), with effectively computable \(C\) and bounded algebraic
degrees & Proved in \href{../TR-BAKER-05.html}{\emph{Effective lower
bounds II: proof of Baker's theorem}}; \(B\) bounds coefficient
heights \\
Feldman's polynomial refinement & Integer-coefficient multiplicative
forms satisfy \(\lvert\Delta\rvert\ge B^{-C}\) & Fixed-base consequence
proved in \href{../TR-BAKER-05.html}{\emph{Effective lower bounds II:
proof of Baker's theorem}}, Section 10; individual-height estimates
treated in the lesson on many logarithms \\
Baker's product-height refinement & An exponent involving
\(\Omega\log\Omega\) & Treated in
\href{../TR-BAKER-08.html}{\emph{Linear forms in many logarithms: the
modern estimates and how to use them}}, alongside the subsequent
refinements \\
Philippon--Waldschmidt, Wüstholz, and Baker--Wüstholz & A lower bound
whose exponent involves \(\Omega\log B\) & Statements, conditions and
deductions given in that same lesson on many logarithms \\
Laurent's determinant method & Two-logarithm estimates involving
\((\log B')^2\) & General parameter theorem proved in
\href{../TR-BAKER-06.html}{\emph{Two logarithms I: an interpolation
determinant}}; sharper explicit families in
\href{../TR-BAKER-07.html}{\emph{Two logarithms II: explicit lower
bounds}} \\
Matveev's estimate & Explicit product-height bounds with improved
dependence on the number of logarithms & Theorem 8.6 and its
applications in the lesson on many logarithms \\
Yu and Bugeaud--Laurent & Upper bounds for a normalized \(p\)-adic
valuation & Proofs planned in \emph{Two p-adic logarithms and Yu's
estimates} \\
\end{longtable}
}

The algebraic-coefficient and integer-coefficient versions use different
size parameters. Keeping this distinction prevents a transcendence
measure from being mistaken for an integer-exponent estimate.

One further distinction matters. For \(\alpha=1+1/N\),

\[
\frac1{N+1}\le\log\left(1+\frac1N\right)\le\frac1N,
\qquad h(\alpha)=\log(N+1).
\]

The inequalities follow by integrating \(1/(1+t)\) from \(0\) to
\(1/N\). The logarithm is small even though the height is large. Refined
estimates with a parameter \(E\) exploit this discrepancy. A height
parameter cannot be replaced by the size of the logarithm alone.

\subsection{7. Exercises}\label{exercises}

\textbf{Exercise 1 (easy).} Determine exactly when \(\Delta=0\) in terms
of \(\Lambda\), and explain what changes for positive real bases. Give
an example showing why principal logarithms of negative bases can
require a period correction.

\textbf{Exercise 2 (easy).} For each \(1\le n\le10\), choose the integer
\(m\ge0\) minimizing \(|2^m-3^n|\). Compare \(|2^m3^{-n}-1|\) with both
elementary rational lower bounds from Section 3.

\textbf{Exercise 3 (medium).} Repeat the box-principle proof for
multiplicatively independent positive rationals
\(\alpha_1,\ldots,\alpha_n\), with \(n\ge2\). Give a sufficient lower
threshold for \(B\) and prove that the nonzero distances from \(1\)
accumulate at zero.

\textbf{Exercise 4 (medium).} In Lemma 1.4, explain why the count is
\((B+1)^n\), why a repeated point is harmless there, and why
multiplicative independence is needed when one wants the constructed
distance to be nonzero. Derive the exact successive-gap bound before
estimating it by \(n\log A/B^{n-1}\).

\textbf{Exercise 5 (hard).} Assume the lower bound in Proposition 1.5.
Prove its explicit search bound, including solutions where one exponent
is zero. Then derive an explicit bound when the hypothesis is weakened
to \(|2^x3^{-y}-1|\ge cB^{-C}\), with known \(0<c\le1\).

\textbf{Exercise 6 (medium).} The bases \(4\) and \(8\) are
multiplicatively dependent. Solve \(4^m-8^n=56\) for integers
\(m,n\ge0\), using the exact power of \(2\) dividing each side. Explain
why the general corollary also applies to this equation.

\subsection{8. Solutions}\label{solutions}

\textbf{Solution 1.} Since \(e^\Lambda=1+\Delta\), the kernel
\(2\pi i\mathbb Z\) of the complex exponential gives the exact
criterion. Real logarithms give a real \(\Lambda\); its intersection
with that kernel is \(\{0\}\). With two bases equal to \(-1\) and both
coefficients equal to \(1\), the principal logarithmic form is
\(2\pi i\), although \(\Delta=0\).

\textbf{Solution 2.} For a fixed \(3^n\), the closest power of \(2\) is
one of the two powers immediately bracketing it. All others are farther
away by monotonicity. Checking those two integers gives:

{\def\LTcaptype{none} % do not increment counter
\begin{longtable}[]{@{}
  >{\raggedright\arraybackslash}p{(\linewidth - 6\tabcolsep) * \real{0.2500}}
  >{\raggedright\arraybackslash}p{(\linewidth - 6\tabcolsep) * \real{0.2500}}
  >{\raggedright\arraybackslash}p{(\linewidth - 6\tabcolsep) * \real{0.2500}}
  >{\raggedright\arraybackslash}p{(\linewidth - 6\tabcolsep) * \real{0.2500}}@{}}
\toprule\noalign{}
\begin{minipage}[b]{\linewidth}\raggedright
\(n\)
\end{minipage} & \begin{minipage}[b]{\linewidth}\raggedright
nearest exponent \(m\)
\end{minipage} & \begin{minipage}[b]{\linewidth}\raggedright
\(\lvert2^m-3^n\rvert\)
\end{minipage} & \begin{minipage}[b]{\linewidth}\raggedright
exact normalized distance
\end{minipage} \\
\midrule\noalign{}
\endhead
\bottomrule\noalign{}
\endlastfoot
1 & 1 or 2 & 1 & \(1/3\) \\
2 & 3 & 1 & \(1/9\) \\
3 & 5 & 5 & \(5/27\) \\
4 & 6 & 17 & \(17/81\) \\
5 & 8 & 13 & \(13/243\) \\
6 & 9 & 217 & \(217/729\) \\
7 & 11 & 139 & \(139/2187\) \\
8 & 13 & 1631 & \(1631/6561\) \\
9 & 14 & 3299 & \(3299/19683\) \\
10 & 16 & 6487 & \(6487/59049\) \\
\end{longtable}
}

Each normalized distance is at least \(3^{-n}\), because its numerator
is a positive integer. It is consequently also at least the weaker bound
\(2^{-m}3^{-n}\). The first two rows attain the actual-denominator
bound. Small coefficient values need not show a smooth trend: the
polynomial lower bounds are uniform restrictions on cancellation, not
monotonicity claims.

\textbf{Solution 3.} Put \(S=\sum_j|\log\alpha_j|>0\). The sums indexed
by \(\{0,\ldots,B\}^n\) lie between
\(B\sum_{\log\alpha_j<0}\log\alpha_j\) and
\(B\sum_{\log\alpha_j>0}\log\alpha_j\), an interval of length \(BS\).
Independence makes the sums distinct. The smallest successive gap is at
most \(BS/((B+1)^n-1)\le S/B^{n-1}\). For an integer \(B>2S\), this is
below \(1/2\). Subtract labels and use Lemma 1.1 to obtain a nonzero
multiplicative distance at most \(2S/B^{n-1}\). These bounds tend to
zero. If the corresponding coefficient vectors belonged to a fixed
finite set, their positive distances would have a positive minimum, a
contradiction. Thus the construction supplies arbitrarily small nonzero
distances.

\textbf{Solution 4.} Each coordinate has \(B+1\) choices, including both
endpoints. There are therefore \((B+1)^n\) labelled sums. A repeated sum
yields a nonzero difference of labels whose product is exactly \(1\),
meeting the weak inequality in the lemma. Under independence, repetition
is impossible. For distinct sums \(s_1<\cdots<s_N\), the \(N-1\)
successive gaps add to \(s_N-s_1\le nB\log A\); at least one gap is at
most \(nB\log A/((B+1)^n-1)\). Finally \((B+1)^n-1\ge B^n\), giving the
stated simpler estimate.

\textbf{Solution 5.} For the original hypothesis, Section 5 gives all
steps. They use only \(x,y\ge0\), so zero exponents are included. With
the weakened hypothesis, the equation gives \(3^y\le(|k|/c)B^C\) and
\(2^x\le(|k|/c)B^C+|k|\le(2|k|/c)B^C\). Replace \(\log(2|k|)\)
throughout the quadratic argument by \(\log(2|k|/c)\). The explicit
bound is

\[
\max(x,y)\le\max\left\{3,
\left(\frac{C+\sqrt{C^2+4\log2\,\log(2|k|/c)}}{2\log2}\right)^2\right\}.
\]

\textbf{Solution 6.} The difference is positive, so \(2m>3n\), and \[
56=2^{3n}(2^{2m-3n}-1).
\] The factor in parentheses is odd. Since \(56=2^3\cdot7\), its exact
power of two forces \(3n=3\), hence \(n=1\). The odd factor then gives
\(2^{2m-3}=8\), so \(m=3\). Conversely, \(4^3-8=56\). This includes the
possible zero-exponent cases: the valuation already forces \(n=1\), and
positivity forces \(m>0\). The corollary requires only a nonzero
difference and a fixed-base lower bound; a multiplicative relation among
the bases does not violate either condition.

\subsection{Prerequisites and
continuation}\label{prerequisites-and-continuation}

The sole arithmetic prerequisite beyond elementary rational arithmetic
is the absolute-height theory specified in Section 3, supplied by the
written \emph{Heights of algebraic numbers} lesson in
\emph{Transcendental numbers}. The local comparison, branch correction,
rational and height lower bounds, box-principle construction,
finite-search calculations and Lemma 1.6 are proved here. The general
difference-of-powers corollary is proved here under its stated
lower-bound hypothesis; the later estimate (5.47) supplies that
hypothesis. The table in Section 6 locates the different theorem shapes
and their conditions. The next two lessons prove the qualitative Baker
theorem before the course turns to effective constants.

\subsection{References}\label{references}

\begin{itemize}
\tightlist
\item
  {[}Waldschmidt 1992{]} Michel Waldschmidt, \emph{Linear Independence
  of Logarithms of Algebraic Numbers}, IMSc Report 116, 1992, revised
  text dated 2016. Introduction and Chapter 3.
  \href{https://webusers.imj-prg.fr/~michel.waldschmidt/articles/pdf/LIL.pdf}{Author's
  text}.
\item
  {[}Waldschmidt 2000{]} Michel Waldschmidt, \emph{Diophantine
  Approximation on Linear Algebraic Groups}, Springer, 2000.
  \href{https://webusers.imj-prg.fr/~michel.waldschmidt/articles/pdf/dalag.pdf}{Author's
  full text}. Section 1.2 explains the elementary denominator,
  box-principle and historical effectivity comparisons; Chapter 3
  supplies the height framework.
\item
  {[}Evertse 2019{]} Jan-Hendrik Evertse, \emph{Diophantine Equations},
  Chapter 5, ``Linear forms in logarithms,'' 2019,
  \href{https://pub.math.leidenuniv.nl/~evertsejh/dio19-5.pdf}{author's
  chapter}. Corollaries 5.3 and 5.5 explain the branch correction and
  normalized difference-of-powers argument. The effective logarithm
  theorem itself is stated there; this course supplies its proof in its
  effective-bound lessons.
\end{itemize}

\subsection{Editable source}\label{editable-source}

\href{TR-BAKER-01.md}{Markdown source}.

\end{document}
